% القسم: المناهج
% الفصل: البراهين
% المبحث: المثال الأول

\documentclass[../../../include/open-logic-section]{subfiles}

\begin{document}

\olfileid{mth}{prf}{ex1}

\olsection{مثال}

مثالنا الأول هو الحقيقة البسيطة الآتية عن اتحادات المجموعات
وتقاطعاتها. وسيوضح لنا كيفية فك التعريفات، وإثبات الاقترانات والقضايا
الكلية، والبرهان بالحالات.

\begin{prop}
لكل مجموعات $A$ و$B$ و$C$،
$A \cup (B \cap C) = (A \cup B) \cap (A \cup C)$.
\end{prop}

فلنثبت ذلك!{}

\begin{proof}
نريد أن نبين أنه لكل مجموعات $A$ و$B$ و$C$،
$A \cup (B \cap C) = (A \cup B) \cap (A \cup C)$.
\begin{quote}
نبدأ بفك تعريف «$=$» في عبارة القضية. تذكر أن إثبات تطابق مجموعتين
يعني بيان أن لهما !!{element} نفسها. أي إن جميع !!{element} المجموعة
$A \cup (B \cap C)$ هي أيضاً !!{element} المجموعة
$(A \cup B) \cap (A \cup C)$، والعكس. ويعني «العكس» أن كل
!!{element} في $(A \cup B) \cap (A \cup C)$ يجب أيضاً أن يكون
!!a{element} في $A \cup (B \cap C)$. وهكذا يتضح، عند فك التعريف، أن
علينا إثبات اقتران. فلنسجل ذلك:
\end{quote}
بحسب التعريف، يكون
$A \cup (B \cap C) = (A \cup B) \cap (A \cup C)$ إذا وفقط إذا كان كل
!!{element} في $A \cup (B \cap C)$ أيضاً !!a{element} في
$(A \cup B) \cap (A \cup C)$، وكان كل !!{element} في
$(A \cup B) \cap (A \cup C)$ أيضاً !!a{element} في
$A \cup (B \cap C)$.
\begin{quote}
بما أن هذا اقتران، فعلينا إثبات كل من طرفيه على حدة. لنبدأ بالأول:
نبرهن أن كل !!{element} في $A \cup (B \cap C)$ هو أيضاً
!!a{element} في $(A \cup B) \cap (A \cup C)$.

هذه قضية كلية، ولذلك نأخذ !!{element} اعتباطياً من
$A \cup (B \cap C)$ ونبين أنه لا بد أن يكون أيضاً !!a{element} في
$(A \cup B) \cap (A \cup C)$. وسنختار متغيراً نسمي به هذا
الـ!!{element} الاعتباطي، وليكن~$z$. ويتواصل برهاننا كما يأتي:
\end{quote}
أولاً، نثبت أن كل !!{element} في $A \cup (B \cap C)$ هو أيضاً
!!a{element} في $(A \cup B) \cap (A \cup C)$. لتكن
$z \in A \cup (B \cap C)$. علينا أن نبين أن
$z \in (A \cup B) \cap (A \cup C)$.
\begin{quote}
حان الآن وقت فك تعريفي $\cup$ و~$\cap$. فتعريف $\cup$، مثلاً، هو:
$A \cup B = \Setabs{z}{z \in A \text{ أو } z \in B}$. وعندما نطبق
التعريف على «$A \cup (B \cap C)$»، يؤدي «$B \cap C$» الآن الدور الذي
كان يؤديه «$B$» في التعريف، ومن ثم
$A \cup (B \cap C) = \Setabs{z}{z \in A \text{ أو }
z \in B \cap C}$. ولذلك يعادل افتراضنا
$z \in A \cup (B \cap C)$ القول:
$z \in \Setabs{z}{z \in A \text{ أو } z \in B \cap C}$.
ولدينا $z \in \Setabs{z}{\dots z\dots}$ إذا وفقط إذا كان
\dots $z$ \dots؛ أي، في هذه الحالة، إذا كان $z \in A$ أو
$z \in B \cap C$.
\end{quote}
بحسب تعريف $\cup$، إما أن $z \in A$ وإما أن $z \in B \cap C$.
\begin{quote}
لما كان هذا فصلاً، فمن المفيد استعمال البرهان بالحالات. نأخذ الحالتين
ونبين أن النتيجة التي نهدف إليها، وهي
«$z \in (A \cup B) \cap (A \cup C)$»، تصح في كل منهما.
\end{quote}
الحالة 1: افترض أن $z \in A$.
\begin{quote}
لا تمنحنا افتراضاتنا مادة أخرى كثيرة. فلننظر إذن إلى ما نحتاج إليه في
النتيجة. نريد أن نبين أن $z \in (A \cup B) \cap (A \cup C)$. وبحسب
تعريف $\cap$، يلزمنا أن نبين أن $z$ ينتمي إلى كل من $(A \cup B)$
و$(A \cup C)$. لكن $z \in A \cup B$ إذا وفقط إذا كان $z \in A$ أو
$z \in B$، ولدينا بالفعل، بوصفه افتراض الحالة~1، أن $z \in A$.
وبالاستدلال نفسه---مع إحلال $C$ محل $B$---نحصل على $z \in A \cup C$.
لقد سرنا بهذا الاستدلال في الاتجاه العكسي، فلنسجله بالاتجاه الذي يحتاج
إليه البرهان.
\end{quote}
بما أن $z \in A$، فإن $z \in A$ أو $z \in B$، ومن ثم، بحسب تعريف
$\cup$، يكون $z \in A \cup B$. وبالمثل $z \in A \cup C$. وهذا يعني،
بحسب تعريف $\cap$، أن $z \in (A \cup B) \cap (A \cup C)$.
\begin{quote}
بهذا تكتمل الحالة الأولى من البرهان بالحالات. ونريد الآن اشتقاق النتيجة
في الحالة الثانية، حيث $z \in B \cap C$.
\end{quote}
الحالة 2: افترض أن $z \in B \cap C$.
\begin{quote}
نتعامل مرة أخرى مع تقاطع مجموعتين. فلنطبق تعريف~$\cap$:
\end{quote}
بما أن $z \in B \cap C$، فلا بد أن يكون $z$ !!a{element} في كل من
$B$ و$C$، بحسب تعريف~$\cap$.
\begin{quote}
حان وقت النظر في النتيجة مرة أخرى. علينا أن نبين أن $z$ ينتمي إلى كل
من $(A \cup B)$ و$(A \cup C)$. والحل مباشر مرة أخرى.
\end{quote}
بما أن $z \in B$، فإن $z \in (A \cup B)$. وبما أن $z \in C$، فإن
$z \in (A \cup C)$ أيضاً. إذن
$z \in (A \cup B) \cap (A \cup C)$.
\begin{quote}
طبقنا هنا تعريفي $\cup$ و$\cap$ مرة أخرى، لكننا سبق أن ذكرناهما وأثبتنا
أن انتماء $z$ إلى إحدى مجموعتين يستلزم انتماءه إلى اتحادهما؛ لذلك لا
نحتاج إلى التصريح بكل ما فعلناه.

لقد أكملنا الحالة الثانية من البرهان بالحالات، فنستطيع الآن تقرير
نتيجتنا الأولى.
\end{quote}
إذن، إذا كان $z \in A \cup (B \cap C)$ فإن
$z \in (A \cup B) \cap (A \cup C)$.
\begin{quote}
لم يبق إلا أن نبين الاتجاه الآخر: كل !!{element} في
$(A \cup B) \cap (A \cup C)$ هو !!a{element} في
$A \cup (B \cap C)$. وكما من قبل، نثبت هذه القضية الكلية بافتراض أن
لدينا !!{element} اعتباطياً من المجموعة الأولى، ثم نبين أنه لا بد أن
ينتمي إلى الثانية. فلنذكر ما سنفعله.
\end{quote}
افترض الآن أن $z \in (A \cup B) \cap (A \cup C)$. نريد أن نبين أن
$z \in A \cup (B \cap C)$.
\begin{quote}
ننطلق الآن من الفرض $z \in (A \cup B) \cap (A \cup C)$. ونأمل ألا
يكون استعمالنا $z$ نفسه هنا كما في الجزء الأول من البرهان مربكاً. فعند
انتهاء ذلك الجزء زال مفعول جميع الافتراضات التي وضعناها فيه، وبوسعنا
الآن وضع افتراضات جديدة عن $z$. وإذا كان هذا مربكاً، فاستبدل بـ$z$
متغيراً آخر فيما يأتي.

نعلم، بحسب تعريف $\cap$، أن $z$ ينتمي إلى كل من $A \cup B$ و
$A \cup C$. وبحسب تعريف $\cup$ نستطيع فك ذلك أكثر فنقول: إما أن
$z \in A$ أو $z \in B$، وكذلك إما أن $z \in A$ أو $z \in C$.
وهذا يبدو كأنه برهان بالحالات مرة أخرى---إلا أن «و» تجعله مربكاً. وقد
تظن أن هذا يعني وجود ثلاثة احتمالات: أن يكون $z$ في $A$ أو $B$ أو
$C$. لكن ذلك خطأ. فعلينا توخي الحذر والنظر في كل فصل على حدة.
\end{quote}
بحسب تعريف $\cap$، لدينا $z \in A \cup B$ و$z \in A \cup C$. وبحسب
تعريف $\cup$، إما أن $z \in A$ أو $z \in B$. ونميز حالتين.
\begin{quote}
بما أننا نركز على الفصل الأول، لم نستخرج بعد فصلنا الثاني، الناتج من
فك $A \cup C$. ولا نحتاج إليه بعد. الحالة الأولى هي $z \in A$، وكل
!!a{element} في مجموعة هو أيضاً !!a{element} في اتحاد تلك المجموعة مع
أي مجموعة أخرى. لذلك فالحالة~1 سهلة:
\end{quote}
الحالة 1: افترض أن $z \in A$. ينتج أن $z \in A \cup (B \cap C)$.
\begin{quote}
والآن إلى الحالة الثانية، $z \in B$. سنفك فيها $\cup$ الثانية ونقيم
برهاناً آخر بالحالات:
\end{quote}
الحالة 2: افترض أن $z \in B$. وبما أن $z \in A \cup C$، فإما أن
$z \in A$ أو $z \in C$. ونميز حالتين فرعيتين:

الحالة 2أ: $z \in A$. وعندئذ، مرة أخرى،
$z \in A \cup (B \cap C)$.
\begin{quote}
حسناً، كان هذا غريباً بعض الشيء. فلم نحتج في هذه الحالة فعلياً إلى
الافتراض $z \in B$، لكن لا بأس بذلك.
\end{quote}
الحالة 2ب: $z \in C$. عندئذ $z \in B$ و$z \in C$، ومن ثم
$z \in B \cap C$، وبالتالي $z \in A \cup (B \cap C)$.
\begin{quote}
بهذا يكتمل البرهانان بالحالات، ونكون قد فرغنا من النصف الثاني.
\end{quote}
إذن، إذا كان $z \in (A \cup B) \cap (A \cup C)$ فإن
$z \in A \cup (B \cap C)$.
\end{proof}

\end{document}
